\scp{Innovations in the morphosyntax}{08-02}
In the morphosyntax \tsc{Taliabu},
\ili{Mangole}, and \ili{Sanana} have preposed genitive proclitics/prefixes
while other languages in the region have postposed enclitics/suffixes.
In \tsc{Taliabu} and \ili{Sanana} the possessed noun takes a prefix
\it{n-} for second or third person possessors (both singular and plural).
This morpheme assimilates to the point of articulation of a following consonant; e.g.
\it{ŋ-kabaŋ} ‘3/2-shield’, \it{ŋ-gege} ‘2/3-armpit’,
\it{m-parabaŋ} ‘2/3-thigh’, \it{m-faha} ‘2/3-arm, hand’ \citep[12--16]{Fortgens1921}.
Examples of preposed and postposed genitive morphology in \tsc{Taliabu}
and \tsc{Sula-Buru} are shown in \trf{08-13},
with third person possessors unless otherwise indicated.

\begin{table}\tcp{Possessive clitics in \tsc{Taliabu} and \tsc{Sula-Buru}}{08-13}
\begin{adjustbox}{width=\textwidth}
	\begin{threeparttable}\st{2pt}
	\begin{tabular}{llllllll}\lsptoprule
local gl.	&	chin, jaw	&	gall, bile	&	grandP	&	father	&	body hair	&	vein	&	bone\\
PMP	&	*qazay	&	*qapəju	&	*upu	&	*ama	&	*bulu	&	*uRat	&	*duRi\\ \midrule
\ili{Soboyo}	&	{\hr}n-ade	&	{\hr}n-oyo-in	&	{\hr}n-uhu	&	{\hr}n-ama	&	{\hr}m-fulu-in	&	{\hr}uha	&	\cg{(talaɲ)}\\
\ili{Taliabu}	&	{\hr}n-ade	&	{\hr}n-oyu-ɲ	&	{\hr}n-uhu	&		&	{\hr}fulu-ɲ	&	{\hr}n-uha-c	&	\cg{(talaŋ)}\\
\ili{Kadai}	&	{\hr}n-ade	&	{\hr}n-oyu	&		&	{\hr}n-ama	&	{\hr}m-fulu	&	{\hr}n-uha	&	\cg{(n-tala)}\\ \midrule
\ili{Mangole}	&		&	{\hr}n-peu	&	{\hr}opu	&		&	{\hr}n-foro foo	&		&	{\hr}loi\\
\ili{Sanana}	&	{\hr}n-aya	&	{\hr}m-peu	&	{\hr}opu	&	{\hr}n-ama	&	{\hr}foa	&	{\hr}ua	&	{\hr}hoi\\
\ili{Lisela}	&	{\hr}aa-n	&	{\hr}peu-n	&	{\hr}upu-n	&	{\hr}ŋ-ama\su{‡}	&	{\hr}ɸulu-n	&		&	{\hr}rohi-n\\
\ili{Buru}	&	{\hr}aa-n	&	{\hr}peu-n	&	{\hr}opo-n	&	{\hr}ama-n	&	{\hr}folo-n	&	{\hr}uha-n	&	{\hr}rohi-n\\ \midrule
\ili{Ambelau}	&	{\hr}ala-mu\su{†}	&	{\hr}feo-ni	&	{\hr}opo-ni	&		&	{\hr}bolu-e	&		&	\cg{(kukui-ni)}\\
		\lspbottomrule
	\end{tabular}
		\begin{tablenotes}\tnsize
			\item[†] {\ili{Ambelau} \it{ala-mu} is `chin/jaw-\tsc{2sg}' = `your chin/jaw'.}
			\item[‡] {\ili{Lisela} \it{ŋ-ama} is borrowed from \ili{Sanana}.}
			%\item[Δ]
			%\item[‖]
			%\item[¶]
		\end{tablenotes}
	\end{threeparttable}
\end{adjustbox}
\end{table}

\citet[68]{Mead2003b} notes that \ili{Banggai} also has the preposed
possessor widely associated with eastern Indonesia,
such as \it{mian lima-no} ‘person hand-\tsc{3sg}’
and \it{koŋgu bonua} ‘my house’.
Other \tsc{\ili{Saluan}-Banggai} languages have post-posed possessors
common to other \tsc{Celebic}
and MP languages in western ISEA.
